
\section{Case Study: Spectral Information Bottleneck (SIB)}
\label{sec:sib}

Section~\ref{sec:evaluation} characterized three failure modes under geographic transfer. We now ask whether this characterization translates into design guidance. We present \sib, a plug-in module that maps each finding to a specific design choice. \sib{} has the smallest mean gap to the within-city upper bound of the ten methods we evaluate, level with FDA. \sib{} is designed from the same LOCO cells on which we evaluate it, so its results show that the diagnosis yields a usable design and do not test the diagnosis independently.

\subsection{Module design}
\label{sec:sib-design}

\begin{figure*}[t]
  \centering
  \includegraphics[width=\linewidth]{fig4_2.png}
  \caption{
\textbf{SIB module and gap to the upper bound.}
\textbf{(a)}~Spectral Information Bottleneck (SIB) plug-in module. Three color-tinted zones group components by the diagnostic finding they target: architecture-aware preprocessing (FDA on CNN inputs, SAG on U-Net skips, both bypassed for ViT); spectral bottleneck on the encoder output (Haar split of features $F$ into content $F_{LL}$ and boundary subbands, content compressed strongly, boundaries gently, $F_{HH}$ passed through); deployment-time coverage threshold $c^*$. Component-to-finding mapping in \S\ref{sec:sib-design}.
\textbf{(b)}~Per-method gap to the within-city upper bound across $9$ leave-one-city-out (LOCO) cells. Open markers: mean gap with $95\%$ CI; filled: worst cell; dashed grey line: Vanilla mean gap. Values in Table~\ref{tab:method-profile}.
FDA: Fourier Domain Adaptation; SAG: Skip Attention Gates.
}
  \label{fig:sib}
\end{figure*}



\textbf{Feature-space failure (\S\ref{sec:feature-space}) $\rightarrow$ band-asymmetric compression.} Based on ~\ref{sec:feature-space}, the intervention must operate on feature space. We need to compress domain-specific content while preserving boundary detail. Haar wavelet decomposition gives this separation directly: the low-frequency subband $F_{LL}$ carries surface appearance and city-specific texture; the high-frequency subbands $F_{LH}, F_{HL}, F_{HH}$ carry shadow geometry. Recent VFM-based domain generalization work confirms this content/style separation in feature space~\cite{fada}. We apply a variational information bottleneck~\cite{alemi2017deep} to $F_{LL}$ aggressively and to $F_{LH}, F_{HL}$ gently. $F_{HH}$ passes through. Aggressive compression on $F_{LL}$ removes city-specific content, while gentle compression on the boundary subbands prevents the bottleneck from erasing the edge structure shadow detection requires.

\textbf{Architecture-dependent locus (\S\ref{sec:architecture-locus}) $\rightarrow$ skip gates and CNN-only preprocessing.} The bottleneck acts on encoder output. In U-Net, skip connections route encoder features directly to the decoder, bypassing the bottleneck. We add Skip Attention Gates: channel-wise sigmoid gates on each skip, identity-initialized so they default to a no-op. The model learns to suppress city-specific channels routed around the bottleneck. CNN encoders also absorb pre-encoder pixel statistics no feature-level intervention can reach. We add Fourier Domain Adaptation~\cite{fda} as input preprocessing on the CNN branch only. ViT encoders carry their recoverable component in the decoder (\S\ref{sec:architecture-locus}) and receive neither component.

\textbf{Class-asymmetric SP failure (\S\ref{sec:sp-gap}) $\rightarrow$ coverage thresholding.} The model's confidence ranking survives geographic shift even when calibration fails. We exploit the surviving ranking with a deployment-time recipe: per cell, fit $c^*_{val}$ as the largest $c$ (coverage) at which selective shadow error on the source-city validation split is at most $(1-\tau) \cdot \text{err}_{val}(c{=}1.0)$, with $\tau = 0.5$ fixed globally, i.e., the criterion is to halve the source-val full-coverage error. Apply $c^*_{val}$ at deployment on the held-out city. The procedure uses no held-out labels. Rank preservation guarantees per-image selective error is bounded by full-coverage error, so the population test measures the magnitude of the recovery rather than its direction. Training loss and full hyperparameters in Appendix~\ref{app:sib-impl}.

\subsection{Results}
\label{sec:sib-results}

\begin{table}[t]
\centering
\small
\begin{minipage}[c]{0.45\textwidth}
\caption{
\textbf{\sib{} component ablations.} Mean $\Delta$ mIoU when each component is removed (or, for the inverse-evidence row, replaced), against each architecture's full-model reference from the same training batch, pooled over the three held-out cities ($n = 450$ paired images). Stars: Holm-adjusted $p$ over the $28$ ablation tests, $^{*}\,p<0.05$, $^{**}\,p<0.01$ (paired image bootstrap, $B = 10{,}000$). $^{\dagger}$MAMNet row removes every VIB. $^{\ddagger}$Run predicts almost no shadow. Full table in Appendix~\ref{app:sib-ablations}.
}
\label{tab:sib-ablations}
\end{minipage}\hfill
\begin{minipage}[c]{0.52\textwidth}
\centering
\setlength{\tabcolsep}{4pt}
\begin{tabular}{l|ccc}
\toprule
\textbf{Variant} & \textbf{MAM} & \textbf{OGL} & \textbf{DIN} \\
\midrule
$-$ SAG                      & $+0.23$ & $-2.06^{**}$ & N/A \\
$-$ FDA, contrast            & $-4.42^{**}$ & $-0.87$ & N/A \\
$-$ Haar                     & $+0.11$ & $-0.63$ & $-23.49^{**\ddagger}$ \\
$-$ VIB on $F_{LL}$          & $+0.29^{\dagger}$ & $-1.01^{**}$ & $-0.71$ \\
VIB on $F_{HL}$ (inv.)       & $-0.29$ & $-1.99^{**}$ & $-0.74$ \\
\bottomrule
\end{tabular}
\end{minipage}
\end{table}

\textbf{Gap to the upper bound.} Over the $9$ cells (3 architectures $\times$ 3 held-out cities; Table~\ref{tab:method-profile}), \sib{} has the smallest mean gap to the within-city upper bound of Vanilla and the eight adaptation methods, $6.42$ mIoU ($95\%$ CI $[5.01, 7.96]$), level with FDA at $6.47$ (difference $+0.05$, $[-0.42, +0.49]$). Its worst-case gap, $10.41$ mIoU (OGLANet, Miami), is second to FDA's $9.10$ and cannot be separated from it (difference $+1.31$, $[-1.01, +3.65]$). FDA and SegDeSiC use unlabeled target-city images in every cell, as \sib{} does on its CNN backbones through FDA. \sib{} leaves a significant gap in all $9$ cells, and on S-EO its worst cell scores $1.57$ mIoU below Vanilla's (Appendix~\ref{app:seo}). The deployment-time coverage recipe reduces selective shadow error on all $9$ cells (mean $\Delta = -0.199$, paired Wilcoxon $p = 0.002$), stable across $\tau \in \{0.3, 0.5, 0.7\}$ (Appendix~\ref{app:tau-sweep}).

\textbf{Component ablations.} Table~\ref{tab:sib-ablations} reports the change in mIoU when each \sib{} component is removed, against each architecture's full-model reference from the same training batch, with Holm correction over the $28$ ablation tests. The largest significant changes are on OGLANet, where removing SAG costs $2.06$ mIoU, compressing $F_{HL}$ in place of $F_{LL}$ costs $1.99$, and removing the content VIB costs $1.01$. On MAMNet the only significant row removes FDA preprocessing together with the contrast stage ($-4.42$ mIoU), and removing SAG changes mIoU by $+0.23$ (not significant). On DINOv3, removing the Haar decomposition gives $-23.49$ mIoU, but that run predicts no shadow on nearly every Chicago and Miami image, so the row records a training collapse and does not isolate the decomposition's contribution; the content-VIB and inverse-subband rows on DINOv3 are not significant after correction. Appendix~\ref{app:sib-ablations} lists every row together with the training differences between each ablation and its reference.

